% P-10/P-11, all-field finite differences and the cubic adjacent pair.
% Dependencies: finding environment, \BR, P-02, R-01, integral convention.
\section{Common-core permanent subspaces and cubic row pencils}

\begin{finding}{P-10}{an integral common-core permanent degeneration}
Let $b\geq2$, $j\geq1$, and let
$e_0,e_1,\ldots,e_{b-2},f_1,\ldots,f_j$ be independent.
Write $Q_i$ for the ordered permanent tensor on
$e_0,e_1,\ldots,e_{b-2},f_i$, and let $a_1,\ldots,a_j$
be independent first-factor vectors. Over every algebraically
closed field,
\[
 \BR\left(\sum_{i=1}^j a_i\otimes Q_i\right)
 \leq (j+1)2^{b-2}.
\]
In particular $\BR(P_b)\leq2^{b-1}$ in every characteristic.
\end{finding}
\begin{proof}
Let $C=\{e_1,\ldots,e_{b-2}\}$, put $S_i=C\cup\{f_i\}$,
and define the tensor-valued polynomial
\[
 F_R(t)=\left(e_0+t\sum_{v\in R}v\right)^{\otimes b},
 \qquad
 D_i(t)=t^{-(b-1)}\sum_{R\subseteq S_i}
                 (-1)^{b-1-|R|}F_R(t).
\]
Inclusion--exclusion kills every monomial missing one
of the $b-1$ indicated directions. Every surviving
monomial has parameter degree at least $b-1$, so the
Laurent expression is polynomial. At degree exactly
$b-1$, every direction occurs once and exactly one
slot contains $e_0$. Its coefficient is therefore
the ordered tensor $Q_i$, with each tensor-basis word
occurring once. Thus $D_i(0)=Q_i$ over $\mathbb Z$.

In $\sum_i a_i\otimes D_i(t)$ group equal evaluation
sets $R$. They are the $2^{b-2}$ subsets of $C$ and,
for each $i$, the $2^{b-2}$ subsets of $C$ with $f_i$
adjoined. Hence there are $(j+1)2^{b-2}$ pure tensor
terms for $t\ne0$, and the limit is the claimed tensor.
No integer is divided out, so the proof survives every
characteristic. The final assertion is the case $j=1$.
\end{proof}

\begin{finding}{P-11}{exact cubic row pencil minima}
For independent $x,y,z,w$ and first-factor basis $a_0,a_1$,
let
\[
 T_+=a_0\otimes\operatorname{per}(x,y,z)
       +a_1\otimes\operatorname{per}(x,y,w),\qquad
 T_-=a_0\otimes\operatorname{alt}(x,y,z)
       +a_1\otimes\operatorname{alt}(x,y,w),
\]
where $\operatorname{per}$ is the sum of all six ordered
words and $\operatorname{alt}$ is their alternating sum.
Then
\[
 \BR(T_+)=6\quad\hbox{in every characteristic},
 \qquad
 \BR(T_-)=8\quad\hbox{if }\operatorname{char}k\ne2.
\]
In characteristic two $T_-=T_+$ and its border rank is six.
Consequently, for every $n\geq4$ the row permanent has
$\mu_2(U_{n,3})=6$; the row determinant has the same
minimum in characteristic two, and minimum eight
in other characteristics. Moreover,
\[
 \BR(P_4)=8\quad\hbox{over every algebraically closed field},
 \qquad \BR(D_4)=8\quad\hbox{in characteristic two}.
\]
\end{finding}
\begin{proof}
The permanent upper bound is P-10 with $b=3,j=2$.
Concretely, the two mixed differences use only the six
evaluation vectors
\[
 x,\quad x+ty,\quad x+tz,\quad x+tw,\quad
 x+ty+tz,\quad x+ty+tw.
\]
Divide the $\{y,z\}$ and $\{y,w\}$ differences of
their third tensor powers by $t^2$, and label them by
$a_0,a_1$. Their constant tensor coefficients are
exactly $T_+$.

For the lower bounds put $A=\langle a_0,a_1\rangle$,
$E=\langle x,y,z,w\rangle$, and use the double-Koszul map
\[
 K_T:A\otimes E\otimes E^*
       \longrightarrow\Lambda^2A\otimes\Lambda^2E\otimes E.
\]
For $T=\sum t_{ijkl}a_i\otimes e_j\otimes e_k\otimes e_l$,
its value on $\alpha\otimes u\otimes\beta$ is
\[
 \sum t_{ijkl}(a_i\wedge\alpha)\otimes(e_j\wedge u)
                   \otimes\beta(e_k)e_l.
\]
Its pure-tensor cost is exactly three: wedge by a nonzero
vector has rank one on the two-dimensional factor and
rank three on the four-dimensional factor, while evaluation
and the remaining fixed vector each have rank one.
Closed determinantal loci therefore give
$\BR(T)\geq\lceil\operatorname{rank}K_T/3\rceil$.

Use zero-based indices, the ordered basis $(x,y,z,w)$,
lexicographically ordered exterior pairs, and product
bases in the displayed factor order. Thus the matrix
has $24$ rows and $32$ columns. For $T_+$ the rows
\[
 (3,8,0,17,9,1,15,7,5,19,11,12,4,13,14,6)
\]
and columns
\[
 (0,1,3,4,5,7,8,9,11,12,13,16,17,20,24,25)
\]
give an order-$16$ minor of determinant $1$.
For $T_-$ the rows
\[
 (3,8,0,17,9,1,15,7,5,19,11,16,12,4,13,2,14,6,18)
\]
and columns
\[
 (0,1,3,4,5,7,8,9,11,12,13,15,16,17,20,21,24,25,28)
\]
give an order-$19$ minor of determinant $4$.
The exact integer matrices and these determinants are
reconstructed from the displayed map by
\path{verify_adjacent_cubic.py}; their index data are
also in \path{adjacent_cubic_certificates.json}.
Thus the permanent lower bound six holds in every
characteristic, and the determinant lower bound seven
holds away from characteristic two.

We now strengthen the determinant lower bound to eight.
The following proper incidence is a necessary condition
for border rank at most seven. Write the three copies of
$E$ as $B,C,D$. There must exist three nine-dimensional
subspaces
\[
 I_{BC}\subset\operatorname{Ann}_{BC}(T_-),\quad
 I_{BD}\subset\operatorname{Ann}_{BD}(T_-),\quad
 I_{CD}\subset\operatorname{Ann}_{CD}(T_-)
\]
of $B^*\otimes C^*$, $B^*\otimes D^*$, and
$C^*\otimes D^*$ respectively, such that
\[
 \dim\bigl(I_{BC}D^*+I_{BD}C^*+I_{CD}B^*\bigr)\leq57.
 \tag{P11.1}
\]
Here products mean insertion into the three tensor
factors, and $\operatorname{Ann}_{BC}(T_-)$ is the
kernel of contraction to $A\otimes D$. To prove the
condition, start with seven generic Segre atoms. Their
pairwise $BC,BD,CD$ projections each span a seven-plane.
The corresponding annihilators are nine-planes, and
their products annihilate all seven independent
$BCD$ projections. Thus their sum has dimension at
most $64-7=57$. The three Grassmannians are projective;
containment in the annihilators and (P11.1) are closed
conditions, the latter being a determinantal condition
on a map from three tautological bundles. Proper
projection therefore extends the condition from a
dense open set of seven-atom sums to its entire secant
closure. This proves the asserted necessary condition
over every algebraically closed field.

Let $H$ be the upper triangular subgroup of
$\operatorname{GL}(E)$ in the basis $(x,y,z,w)$.
It preserves the plane
$U=\langle\operatorname{alt}(x,y,z),
\operatorname{alt}(x,y,w)\rangle$. Give $A=U^*$ its
contragredient action. Then $H$ fixes $T_-$ and
preserves the projective incidence fiber just constructed.
This fiber has an $H$-fixed point, as follows directly
from triangular matrices. Embed the product of the
three Grassmannians by its Segre--Pl\"ucker embedding.
The ambient representation is a tensor product of
exterior powers of triangular representations, hence
has an $H$-stable complete flag $0=V_0\subset\cdots
\subset V_N$. Choose the least $s$ for which the fiber
meets $\mathbb P(V_s)$. This intersection is projective
and lies in the affine chart
$\mathbb P(V_s)\setminus\mathbb P(V_{s-1})$, so it is
finite. Since $H$ is connected, each of its points is
fixed. Thus all three nine-planes may simultaneously
be assumed $H$-stable. This argument uses triangularity
and properness in every characteristic; it makes no
semisimplicity assumption.

Let $\epsilon_0,\epsilon_1,\epsilon_2,\epsilon_3$ be
the dual basis. For any pair of the three $E$ factors put
\[
 S_{ii}=\epsilon_i\otimes\epsilon_i,
 \qquad
 S_{ij}=\epsilon_i\otimes\epsilon_j+
         \epsilon_j\otimes\epsilon_i\quad(i<j),
 \qquad
 Z=\epsilon_2\otimes\epsilon_3-
       \epsilon_3\otimes\epsilon_2.
\]
When the characteristic is not two, direct contraction
gives the eleven-dimensional space
\[
 \operatorname{Ann}_{BC}(T_-)
   =\langle S_{ij}:0\leq i\leq j\leq3\rangle\oplus kZ,
 \tag{P11.2}
\]
and the same description for the other two pairs.
Indeed the pair contraction has the five independent
images $a_0\otimes x,a_0\otimes y,a_1\otimes x,
a_1\otimes y,a_0\otimes z+a_1\otimes w$.

There are exactly two $H$-stable nine-planes in (P11.2):
\[
 \begin{split}
 J_0&=\langle S_{ij}:(i,j)\ne(0,0)\rangle,\\
 J_1&=\langle S_{ij}:(i,j)\notin\{(0,0),(0,1)\}\rangle
          \oplus kZ.
 \end{split}
 \tag{P11.3}
\]
Here is the full classification. The diagonal torus
has a one-dimensional weight space at each $S_{ij}$,
except that the $(2,3)$ weight also contains $Z$.
An invariant nine-plane cannot contain $S_{00}$:
its upper-triangular orbit spans all ten symmetric
vectors. If it contains $S_{01}$, that orbit spans
all nine symmetric vectors except $S_{00}$, so the
plane equals $J_0$. To check the latter statement,
the elementary substitution
$\epsilon_0\mapsto\epsilon_0+t\epsilon_1$ gives
$S_{01}+2tS_{11}$, and substitutions with larger
indices generate the other eight vectors. Two is
invertible here. If the plane contains neither
$S_{00}$ nor $S_{01}$, torus invariance excludes both
of those one-dimensional weight spaces. The remaining
space has dimension nine, forcing the plane to be
$J_1$. Both displayed spaces are stable, which also
verifies that the list is exact.

It remains to test the eight possible triples from
(P11.3). The multiplication map in (P11.1) has $64$
rows and $108$ columns. Order the rows lexicographically
in $B^*,C^*,D^*$; order the columns by pair $BC,BD,CD$,
then the basis in (P11.3) with the $S_{ij}$ in
lexicographic order and $Z$ last, then the free index
$0,1,2,3$. The complete certificate supplies an
order-$58$ minor for each triple, as follows:
\[
\begin{array}{c|rrrrrrrr}
 (J_{BC},J_{BD},J_{CD})
 &(0,0,0)&(0,0,1)&(0,1,0)&(0,1,1)
 &(1,0,0)&(1,0,1)&(1,1,0)&(1,1,1)\\ \hline
 \det&-4&4&-4&8&16&32&-64&-64
\end{array}
\]
The rows and columns are recorded in
\path{adjacent_cubic_borel_certificate.json}.
The standalone \path{verify_adjacent_cubic.py}
reconstructs every entry directly from (P11.3) and
evaluates these integer determinants by fraction-free
elimination, without importing the generator or using
a rank oracle. Every displayed minor is nonzero when
two is invertible, giving rank at least $58$ in all
eight cases and contradicting (P11.1). Hence
$\BR(T_-)\geq8$ in every odd characteristic and in
characteristic zero.

For the determinant upper bound put $z_0=z,z_1=w$ and
consider the two simultaneous tangent tensors
\[
 H_x=\sum_{i=0}^1a_i\otimes
 (z_i\otimes x\otimes y-x\otimes z_i\otimes y
                         +x\otimes x\otimes z_i),
\]
\[
 H_y=\sum_{i=0}^1a_i\otimes
 (-z_i\otimes y\otimes x+y\otimes z_i\otimes x
                         -y\otimes y\otimes z_i).
\]
Each has border rank at most three. For $H_x$, the
three-term curve is
\[
 \frac{1}{t}\left[
 \sum_{i=0}^1a_i\otimes(x+tz_i)\otimes(x-tz_i)\otimes(y+tz_i)
 -(a_0+a_1)\otimes x\otimes x\otimes y\right].
\]
For $H_y$, it is
\[
 \frac{1}{t}\left[
 \sum_{i=0}^1a_i\otimes(y-tz_i)\otimes(y+tz_i)\otimes(x-tz_i)
 -(a_0+a_1)\otimes y\otimes y\otimes x\right].
\]
Both numerators have zero constant coefficient, so
the displayed expressions extend polynomially over
$t=0$, with limits $H_x,H_y$. Direct expansion gives
the integral identity
\[
 T_-=H_x+H_y-
       \sum_{i=0}^1a_i\otimes(x+y)\otimes(x-y)\otimes z_i.
\]
The last sum consists of two pure tensors. Adding it
to the two curves produces an eight-term degeneration
of $T_-$. This upper bound in fact holds in every
characteristic: only the parameter $t$ is inverted,
and no integer is divided out. The standalone checker
reconstructs all eight terms, verifies cancellation of
every negative power, and checks all $128$ constant
tensor coefficients over $\mathbb Z$.

P-02 identifies every two-dimensional row minimum with
this adjacent pair. For $P_4$, R-01 with $b=3,j=2$
gives $\BR(P_4)\geq4-2+6=8$; P-10 with $b=4,j=1$
gives the matching upper bound eight. Finally the two
row tensors agree in characteristic two.
\end{proof}

\paragraph{Geometric dictionary.}
P-10 is an integral family of Segre sections over
$\mathbb G_m$ whose inclusion--exclusion sum extends
over $\mathbb A^1$. P-11 uses exterior-power maps of
locally free sheaves and their actual integer minors;
both upper degenerations are integral polynomial
families after their parameter poles cancel. The exact
foundations are EGA~I \S9.5,
SGA~3 Expos\'e~I \S\S4.5--4.6, and SGA~6 Expos\'e~V,
Definition~2.1, Examples~2.2(b) and Lemma~2.2.1, with
the printed-page and facsimile locators in the main
dictionary. The rank-eight row corollary further uses
R-01's EGA~II proper incidence. The determinant lower
bound uses a projective product of Grassmannians,
tautological multiplication maps, and a triangular
Segre--Pl\"ucker representation; the finite incidence
projection is proper by EGA~II, Proposition~5.4.2 and
Theorem~5.5.3(i). Its eight power-of-two minors work
over $\mathbb Z[1/2]$. No archive cote or
cohomological rank theorem is inferred from these
constructions.
